% ------------------------------------------------------------------------
% abnTeX2: Modelo de Projeto de pesquisa em conformidade com
% ABNT NBR 15287:2011 Informação e documentação - Projeto de pesquisa -
% Apresentação
% ------------------------------------------------------------------------
\documentclass[
	12pt,           % tamanho da fonte
	openright,      % capítulos começam em pág ímpar
	oneside,        % para impressão em recto (frente) e verso (atrás) -> nesse caso, "oneside"
	a4paper,        % tamanho do papel
	brazil          % o último idioma é o principal do documento
]{abntex2}

% ----------------------------------------------------------
% PACOTES
% ----------------------------------------------------------
\usepackage{lmodern}            % Usa a fonte Latin Modern
\usepackage[T1]{fontenc}        % Seleção de códigos de fonte
\usepackage[utf8]{inputenc}     % Codificação do documento (conversão automática dos acentos)
\usepackage{indentfirst}        % Indenta o primeiro parágrafo de cada seção
\usepackage{color}              % Controle das cores
\usepackage{graphicx}           % Inclusão de gráficos
\usepackage{microtype}          % Para melhorias de justificação
\usepackage{lipsum}             % Para geração de dummy text (exemplo)
\usepackage[brazilian,hyperpageref]{backref}	
\usepackage[alf]{abntex2cite}	% Citações padrão ABNT
\usepackage{hyperref}           % Hyperlinks no PDF
\usepackage{multirow}           % Para tabelas com células que ocupam várias linhas

% ----------------------------------------------------------
% CONFIGURAÇÕES DE PACOTES
% ----------------------------------------------------------

% BACKREF
\renewcommand{\backrefpagesname}{Citado na(s) página(s):~}
\renewcommand{\backref}{}
\renewcommand*{\backrefalt}[4]{
  \ifcase #1
    Nenhuma citação no texto.
  \or
    Citado na página #2.
  \else
    Citado #1 vezes nas páginas #2.
  \fi
}
\renewcommand\thesection{\arabic{section}}

% HYPERSETUP
\hypersetup{
	pdftitle={Projeto de Iniciação Científica: Dashboard para Monitoramento de Fábricas de Software Acadêmicas},
	pdfauthor={Danrley Pereira e Kerlla Luz},
	pdfsubject={Projeto de pesquisa submetido ao edital de iniciação científica do UDF},
	pdfkeywords={Fábrica de Software Acadêmica, Desenvolvimento de software, Educação em TI, Dashboard, Métricas},
	colorlinks=true,
	linkcolor=blue,
	citecolor=blue,
	filecolor=magenta,
	urlcolor=blue
}

% ---
% Informações de dados para CAPA e FOLHA DE ROSTO
% ---
\titulo{Desenvolvimento de um Dashboard para o Monitoramento e Acompanhamento Contínuo de Projetos em Fábricas de Software Acadêmicas}
\autor{Danrley Willyan da Silva Pereira \and Kerlla Luz}
\local{Brasília}
\data{2025}
\instituicao{
  Centro Universitário do Distrito Federal (UDF)\\
  Curso de Ciência da Computação\\
  Programa de Iniciação Científica
}
\tipotrabalho{Projeto de Pesquisa}
\preambulo{Projeto de pesquisa submetido ao Edital do Programa Institucional de Bolsas de Iniciação Tecnológica \& Inovação de 2025 do Centro Universitário do Distrito Federal (UDF).}

% ---
% Início do Documento
% ---
\begin{document}
\selectlanguage{brazil}
\frenchspacing

% ----------------------------------------------------------
% ELEMENTOS PRÉ-TEXTUAIS
% ----------------------------------------------------------
\imprimircapa
\imprimirfolhaderosto

% ----------------------------------------------------------
% SUMÁRIO
% ----------------------------------------------------------
\pdfbookmark[0]{\contentsname}{toc}
\tableofcontents*
\cleardoublepage

% ----------------------------------------------------------
% ELEMENTOS TEXTUAIS
% ----------------------------------------------------------
\textual

% ----------------------------------------------------------
% INTRODUÇÃO
% ----------------------------------------------------------
\section[Introdução]{Introdução}
\addcontentsline{toc}{chapter}{Introdução}

O setor de Tecnologia da Informação (TI) tem desempenhado um papel cada vez mais relevante no desenvolvimento econômico e social, especialmente em polos emergentes como Brasília \cite{ibge2025}, que abriga diversas instituições públicas, privadas e centros de pesquisa. Nesse contexto, a criação de uma \textbf{Fábrica de Software Acadêmica} em ambiente universitário surge como oportunidade estratégica para integrar teoria e prática, preparando os estudantes para os desafios reais do mercado ao mesmo tempo em que se desenvolvem soluções voltadas a demandas concretas da comunidade \cite{dos2021experiencia, neto2017relato}.

Instituições de ensino superior identificam a falta de experiência prática dos alunos como um gap a ser preenchido \cite{romanha2016modelo}. A implementação de Fábricas de Software Acadêmicas surge da necessidade de proporcionar experiência prática aos alunos, o que está alinhado às diretrizes curriculares do MEC \cite{dcncomputacao2016}; assim, as instituições que melhor oferecem um ambiente de aprendizado prático tendem a cumprir satisfatoriamente a missão de preparar bons profissionais para o mercado de trabalho. Em outras palavras, criar um ambiente de “quase indústria” dentro da universidade ajuda a formar graduandos mais preparados para os desafios reais da profissão.

Muitas iniciativas acadêmicas adotam projetos de Fábrica de Software exatamente “com o intuito de fornecer aos alunos experiências do mercado de trabalho durante a formação universitária” \cite{dos2021experiencia}. Esses projetos integradores simulam ambientes profissionais, permitindo que os estudantes apliquem conhecimentos teóricos em projetos de software reais. Um exemplo foi a criação de um programa de “residência em desenvolvimento de software” em que o aluno vivencia “uma experiência profissional em um ambiente próximo do encontrado no mercado de trabalho” \cite{dos2021experiencia}. Tais ambientes acadêmicos simulados expõem os alunos a práticas, ferramentas e ritmos de trabalho típicos da indústria de software, encurtando a distância entre a sala de aula e o mundo profissional.

Dentre essas iniciativas, existe o projeto da \emph{Fábrica de Software Acadêmica (LabTech)} do Centro Universitário do Distrito Federal (UDF), projeto qual tem evidenciado a importância de disponibilizar métricas de progresso, colaboração e produtividade tanto para professores quanto para gerentes e integrantes das equipes. Enfatizando esse ponto, os relatórios de alunos e professores deixam claro a ausência de um registro centralizado e de visualizações claras, o que dificulta a tomada de decisão e análises do progresso dos projetos.

Nesse sentido, o desenvolvimento de metodologias de coleta e análise de dados que suportem o processo de desenvolvimento de software tem sido alvo de constantes pesquisas em Fábricas de Software de Instituições de Ensino Superior (IES) no Brasil \cite{romanha2016modelo, neto2017relato}. Além de tornar o aprendizado mais significativo para os estudantes, a adoção de métricas claras pode contribuir para a transparência de resultados e a melhoria contínua dos projetos \cite{oliveira2017conduccao, almeida2006liccoes}.

Assim, neste projeto de Iniciação Tecnológica, propõe-se o desenvolvimento de um \textbf{dashboard} que colete dados diários (ou semanais) do GitHub\footnote{GitHub. Disponível em: \url{https://github.com/about}. Acessado em: 26 de fevereiro de 2025.} (repositório de código fonte, ver \ref{sec:tecnico}), a fim de exibir \textbf{métricas de colaboração e performance} de forma integrada. Essa solução visa beneficiar professores, gestores e os próprios estudantes, permitindo melhor acompanhamento do desenvolvimento dos projetos e fomentando práticas de melhoria contínua nas Fábricas de Software Acadêmicas.

\subsection{Motiva\c{c}\~{a}o}
O aluno proponente junto com a orientadora têm um papel ativo no LabTech, sendo que o aluno proponente exerceu o papel de gerente no primeiro semestre de 2023 e em ambos semestres de 2024. Com base nessa experi\^{e}cia, ambos identificaram gargalos no acompanhamento dos alunos: como falta de m\'{e}tricas de desempenho, falta de análises em tempo real do progresso do projeto e das atividades, dificuldade de uso das ferramentas por parte dos integrantes da fábrica, entre outras dificuldades. Isso exposto, a disponibilidade de um dashboard com métricas em tempo real pode auxiliar no acompanhamento dos alunos, criar planos de ação, buscar nivelar o nível dos alunos, identificar silos de conhecimento etc.

\subsection*{Objetivo Geral}
Criar uma \textbf{solução tecnológica (dashboard)} para monitoramento e acompanhamento contínuo de projetos de software acadêmicos, possibilitando a visualização de métricas de progresso, capacidade e engajamento das equipes, a fim de potencializar a gestão e o aprendizado em Fábricas de Software Acadêmicas.

\subsubsection*{Objetivos Específicos}
\begin{enumerate}
    \item \textbf{Realizar uma revisão de literatura} sobre métodos e métricas adotados em projetos de desenvolvimento de software, com foco nas Fábricas de Software Acadêmicas, a fim de identificar boas práticas para coleta e interpretação de dados e embasar a definição das métricas que serão apresentadas no dashboard.
    \item \textbf{Definir o conjunto de indicadores-chave (KPIs)} que permitam avaliar tanto a parte técnica (\emph{commits}, \emph{issues}, \emph{PRs}, bugs) quanto a parte colaborativa (rede de colaboração, frequência de comunicação).
    \item \textbf{Implementar um protótipo inicial} de coleta automática de dados a partir da API do GitHub, desenvolvendo \emph{scripts} ou serviços que atualizem um banco de dados central.
    \item \textbf{Construir o dashboard} que apresente as informações em tempo real (ou quase real), classificando os resultados por projeto, por equipe e por indivíduo, de maneira que gestores e professores possam identificar rapidamente gargalos e boas práticas.
    \item \textbf{Conduzir testes-piloto} com equipes participantes da Fábrica de Software do UDF (ou de outra instituição parceira), documentando ajustes técnicos e sugerindo melhorias no modelo de visualização de dados.
\end{enumerate}

Além dessas metas, este projeto prevê a análise dos resultados, a elaboração de um relatório final, a disponibilização do produto sob \textbf{licença GPL3} e a apresentação dos achados no congresso científico do UDF.

\subsection{Dashboard Open Source sob Licença GPL 3: Flexibilidade e Extensibilidade}

Uma das principais vantagens da nossa proposta é que o dashboard será desenvolvido como uma solução open source sob a licença GNU General Public License versão 3 (GPL 3). A GPL 3 é uma licença copyleft que garante a liberdade de uso, estudo, modificação e distribuição do software, desde que todas as versões derivadas mantenham as mesmas liberdades. Essa característica permite que o código seja estendido e adaptado para atender a diferentes contextos, possibilitando que qualquer fábrica ou equipe de software no mundo possa colocar a ferramenta em produção sem incorrer em custos de licenciamento.

\subsection{Uso de Dashboards para Gestão de Projetos Acadêmicos: Benchmark de Ferramentas}

A literatura e o mercado atual oferecem uma variedade de soluções que integram métricas extraídas de repositórios—como GitHub, GitLab e Bitbucket—em dashboards dinâmicos, permitindo a visualização em tempo real do progresso dos projetos e a identificação de gargalos. Tais soluções são fundamentais para reduzir a assimetria de informação e promover a colaboração entre as equipes.

Entre as soluções SaaS disponíveis, destaca-se o \textbf{Datadog}\footnote{Datadog. Disponível em: \url{https://www.datadoghq.com/}. Acessado em: 26 de fevereiro de 2025.}, que integra dados de múltiplas fontes (incluindo repositórios) para oferecer uma visão abrangente dos processos de desenvolvimento. De forma similar, o \textbf{New Relic}\footnote{New Relic. Disponível em: \url{https://newrelic.com/}. Acessado em: 26 de fevereiro de 2025.} proporciona monitoramento detalhado e uma interface amigável para acompanhar indicadores-chave, como tempo de ciclo e taxa de resolução de \emph{issues}.

Plataformas como o \textbf{Keen IO}\footnote{Keen IO. Disponível em: \url{https://keen.io/}. Acessado em: 26 de fevereiro de 2025.} demonstram a eficácia de dashboards customizáveis, que possibilitam a consolidação de métricas técnicas (por exemplo, número de \emph{commits} e \emph{pull requests}) e colaborativas (como a frequência de interações em \emph{issues}). Outras soluções focadas em repositórios, como o \textbf{iftrue.co}\footnote{iftrue.co. Disponível em: \url{https://www.iftrue.co/}. Acessado em: 26 de fevereiro de 2025.} e o \textbf{ScatterSpoke Metrics for Bitbucket}\footnote{ScatterSpoke Metrics for Bitbucket. Disponível em: \url{https://www.scatterspoke.com/}. Acessado em: 26 de fevereiro de 2025.}, utilizam análises baseadas em inteligência artificial para identificar gargalos e otimizar os fluxos de trabalho.

Adicionalmente, ferramentas que agregam dados de múltiplas fontes e plataformas voltadas à análise integrada de métricas, como o \textbf{Worklytics}\footnote{Worklytics. Disponível em: \url{https://www.worklytics.co/}. Acessado em: 26 de fevereiro de 2025.}, demonstram que a consolidação e a visualização clara de indicadores contribuem para uma gestão mais transparente e colaborativa dos projetos.

\subsection{Benchmark de Soluções: Comparação de Preços e Funcionalidades}

Para demonstrar a importância do desenvolvimento de um dashboard open source, é útil comparar as soluções disponíveis no mercado em termos de funcionalidades e custos. A Tabela~\ref{tab:benchmark} apresenta um resumo dos preços e principais recursos de algumas ferramentas SaaS que agregam métricas de repositórios (como GitHub, GitLab e Bitbucket), bem como alternativas que visam oferecer dashboards interativos.

\begin{table}[htbp]
\centering
\caption{Comparativo de Soluções de Dashboard para Monitoramento de Repositórios}
\label{tab:benchmark}
\begin{tabular}{|p{3cm}||p{3cm}|p{3cm}|p{2cm}|}
\hline
\textbf{Ferramenta} & \textbf{Tipo} & \textbf{Principais Funcionalidades} & \textbf{Preço (a partir de)} \\ \hline
Datadog\footnote{Datadog. Disponível em: \url{https://www.datadoghq.com/}. Acessado em: 26 de fevereiro de 2025.}  
  & SaaS & Monitoramento em tempo real, integração com múltiplas fontes (incl. Bitbucket, GitHub) & \$15/host/mês \\ \hline
New Relic\footnote{New Relic. Disponível em: \url{https://newrelic.com/}. Acessado em: 26 de fevereiro de 2025.} 
  & SaaS & APM, monitoramento de infraestrutura, dashboards customizáveis & Plano gratuito + cobrança por uso \\ \hline
Keen IO\footnote{Keen IO. Disponível em: \url{https://keen.io/}. Acessado em: 26 de fevereiro de 2025.} 
  & SaaS & APIs flexíveis para coleta e visualização customizada de dados & \$149/mês \\ \hline
iftrue.co\footnote{iftrue.co. Disponível em: \url{https://www.iftrue.co/}. Acessado em: 26 de fevereiro de 2025.} 
  & SaaS & Métricas de experiência do desenvolvedor, análises de DORA & Gratuito para até 10 desenvolvedores; \$90/team/mês \\ \hline
ScatterSpoke Metrics\footnote{ScatterSpoke Metrics for Bitbucket. Disponível em: \url{https://www.scatterspoke.com/}. Acessado em: 26 de fevereiro de 2025.} 
  & SaaS & Insights baseados em IA para otimização de workflow & A partir de \$90/team/mês \\ \hline
Worklytics\footnote{Worklytics. Disponível em: \url{https://www.worklytics.co/}. Acessado em: 26 de fevereiro de 2025.} 
  & SaaS & Análise de produtividade e bem-estar da equipe & Preço não divulgado \\ \hline
\end{tabular}
\end{table}

\noindent
\textbf{Discussão:}  
Como pode ser observado, soluções comerciais como \emph{Datadog} e \emph{New Relic} oferecem dashboards robustos com monitoramento em tempo real e ampla integração com diversas plataformas. No entanto, os custos—que podem variar desde \$15 por host/mês até modelos de cobrança baseados no volume de dados—podem ser um entrave para equipes com orçamentos limitados, especialmente em contextos acadêmicos. Por outro lado, soluções SaaS como \emph{Keen IO}, \emph{iftrue.co} e \emph{ScatterSpoke Metrics} apresentam preços mais definidos (a partir de \$90 a \$149/mês por equipe) e focam na customização e na análise detalhada dos dados de repositórios.

Diante desse cenário, o desenvolvimento de um \textbf{dashboard open source} sob licença GPL 3 ganha ainda mais relevância, pois pode oferecer funcionalidades comparáveis—como a consolidação de indicadores técnicos (commits, issues, pull requests) e colaborativos (frequência de comunicação, rede de colaboração)—sem os custos recorrentes das soluções comerciais. Essa abordagem não só reduz o investimento financeiro, como também permite maior flexibilidade para adaptações específicas ao contexto das Fábricas de Software Acadêmicas, promovendo a transparência e a melhoria contínua dos projetos.




% ----------------------------------------------------------
% CAPÍTULO: REVISÃO DE LITERATURA
% ----------------------------------------------------------
\section{Revisão de Literatura}

\subsection{Fábricas de Software Acadêmicas}
Fábricas de Software têm sido propostas como ambientes de desenvolvimento de projetos reais em ambiente acadêmico, unindo teoria e prática para o aprimoramento de estudantes. No contexto brasileiro, diversos estudos destacam suas vantagens na formação de profissionais de TI, pois permitem experimentar metodologias ágeis, adoção de ferramentas modernas e uso de \emph{frameworks} \cite{oliveira2017conduccao, fernandes2004fabrica, romanha2015fabrica, almeida2006liccoes}.

Além disso, vários autores salientam que essas iniciativas aproximam os estudantes das demandas do mercado e da comunidade, ao mesmo tempo em que promovem um aprendizado ativo e colaborativo \cite{dos2021experiencia, neto2017relato}. Assim, projetos de Fábricas de Software Acadêmicas podem oferecer uma experiência mais realista sobre como ocorre o ciclo de vida de projetos de software, contribuindo para a formação de competências alinhadas às exigências do setor de TI \cite{romanha2016modelo}.

\subsection{Métricas em Projetos de Software Acadêmicos}
Apesar das vantagens mencionadas, a falta de padronização nas práticas de medição e acompanhamento ainda é considerada um desafio. \citeonline{kitchenham2010s} destaca a importância de adotar métodos científicos de coleta e análise de dados para embasar a evolução de projetos dentro das Fábricas. Nesse sentido, estabelecer \emph{KPIs} que reflitam tanto o aspecto técnico (ex.: \emph{commits}, \emph{issues}, \emph{bugs}, \emph{merges}, tempo de revisão) \cite{hamer2021using, barbosaIndicadoresConjunto2016}, quanto o aspecto colaborativo (rede de colaboração, frequência de comunicação em \emph{pull requests} e \emph{issues}, etc.) pode oferecer um panorama claro do andamento de cada projeto \cite{basili1996validation}.

Segundo \citeonline{wef:2025}, as principais \emph{soft skills} requeridas no mercado de TI incluem pensamento analítico, resiliência, liderança e colaboração. Projetos educacionais podem se beneficiar de métricas que avaliem essas competências de modo indireto, como o número de revisões de código realizadas, a distribuição das tarefas por membro e a constância de participação nas discussões técnicas.

Ainda nessa linha, \citeonline{hamer2021using} e \citeonline{barbosaIndicadoresConjunto2016} ressaltam que mensurar elementos tanto quantitativos quanto qualitativos (por exemplo, a participação em code reviews, o fechamento de \emph{issues}, a comunicação entre membros) permite um retorno construtivo aos estudantes. Essa retroalimentação subsidia \emph{feedbacks} contínuos e serve de base para ajustes pedagógicos \cite{almeida2006liccoes}.

\subsection{Uso de Dashboards para Gestão de Projetos Acadêmicos}
A literatura sobre ferramentas de apoio à gestão de projetos, como \emph{Jira}, \emph{Trello} e dashboards personalizados, mostra que uma visualização clara de indicadores reduz a assimetria de informação, facilita a detecção de gargalos e incentiva a colaboração \cite{temitope2020software, dos2020usoTrello, ozkan2019agile}. Em âmbito acadêmico, porém, faltam estudos que avaliem a eficácia de dashboards voltados especificamente para análise de engajamento e produtividade em Fábricas de Software Acadêmicas.

O desenvolvimento de \emph{dashboards} que consolidem as diversas métricas de forma interativa pode ajudar docentes, coordenadores e gerentes de Fábricas Acadêmicas a tomar decisões mais embasadas. Ao correlacionar dados de \emph{commits}, \emph{issues} e participação em discussões técnicas, os responsáveis podem identificar necessidades de intervenção, distribuição ineficiente de tarefas ou lacunas de conhecimento na equipe \cite{kitchenham2010s}.

\subsection{Detalhes técnicos}
\label{sec:tecnico}

Existe uma variedade de ferramentas de hospedagem de repositórios de código-fonte disponíveis na web. Entretanto, é fundamental compreender que os serviços de hospedagem e os sistemas de controle de versão são entidades distintas. Os sistemas de controle de versão são os utilitários utilizados para gerenciar as alterações ocorridas ao longo do ciclo de desenvolvimento de software, registrando as modificações efetuadas nos arquivos de código-fonte – exemplos clássicos são o Git\footnote{Git. Disponível em: \url{https://git-scm.com/}. Acessado em 26 de fevereiro de 2025.} e o SVN\footnote{SVN. Disponível em: \url{https://subversion.apache.org/}. Acessado em: 26 de fevereiro de 2025.}. Por outro lado, os serviços de hospedagem de repositórios são aplicações web que integram e potencializam as funcionalidades dos sistemas de controle de versão, oferecendo recursos adicionais para colaboração e visualização de indicadores, como o GitLab\footnote{GitLab. Disponível em: \url{https://about.gitlab.com/}. Acessado em: 26 de fevereiro de 2025.} e o GitHub. No âmbito desse projeto, o GitHub foi escolhido como repositório para extração de dados, em virtude de sua ampla popularidade na comunidade de desenvolvedores e por ser usado internamente no LabTech.

% ----------------------------------------------------------
% CAPÍTULO: METODOLOGIA
% ----------------------------------------------------------
\section{Metodologia}

A metodologia adotada neste projeto segue uma abordagem mista, considerando tanto aspectos qualitativos (observações, relatórios, diários de bordo) quanto quantitativos (métricas coletadas automaticamente do GitHub).  

\subsection{Tipo de Pesquisa}
O trabalho caracteriza-se como pesquisa aplicada, com foco em produzir uma ferramenta (dashboard) de uso prático em Fábricas de Software. Será ainda de natureza \textbf{descritiva} e \textbf{exploratória}, pois busca mapear e sintetizar métricas relevantes de sucesso em projetos acadêmicos.

\subsection{Coleta de Dados}
A coleta de dados será realizada em duas vertentes:
\begin{enumerate}
    \item \textbf{Qualitativa:} Entrevistas breves com professores e gestores, além de observações em reuniões do LabTech, visando entender quais métricas são mais valorizadas e quais necessidades não estão sendo atendidas por ferramentas existentes.
    \item \textbf{Quantitativa:} Uso da \emph{API do GitHub} para extrair dados de \emph{commits}, \emph{issues}, \emph{pull requests}, comentários e dados de rede (quem revisa o código de quem, quem abre \emph{issues} para qual membro etc.).
\end{enumerate}

\subsection{Análise e Desenvolvimento}
\begin{itemize}
    \item \textbf{Definição de KPIs}: A partir da literatura e das entrevistas, serão selecionadas as principais métricas que melhor reflitam progresso (ex.: burndown, velocidade de fechamento de \emph{issues}), engajamento (rede de colaboração, distribuição de \emph{commits}) e aspectos qualitativos.
    \item \textbf{Prototipagem}: Desenvolvimento de \emph{scripts} em \textit{Python} ou \textit{JavaScript} para coleta e consolidação dos dados em um banco de dados. Posteriormente, criação de uma interface web (ex.: \emph{React} ou \emph{Vue.js}) para visualização das métricas.
    \item \textbf{Testes-Piloto}: Implementação inicial do dashboard para uso pelo LabTech, coletando feedback dos participantes e ajustando funcionalidades e visualizações.
\end{itemize}

\subsection{Aspectos Éticos}
Será garantido o anonimato dos participantes nas análises de engajamento (por exemplo, usando \emph{hash} de e-mails ou identificadores). Os dados coletados serão utilizados apenas para fins de pesquisa e melhoria de processos, sem divulgação pública de informações sensíveis. O projeto segue as diretrizes éticas recomendadas pelo Comitê de Ética em Pesquisa do UDF.

% ----------------------------------------------------------
% CAPÍTULO: CRONOGRAMA E RECURSOS
% ----------------------------------------------------------
\section{Cronograma e Recursos}

O período previsto para execução do projeto inicia-se em abril de 2025 (mês 1) e encerra-se em março de 2026 (mês 12), totalizando 12 meses. Abaixo descrevemos a sequência de atividades:

\subsection{Recursos Necessários}
\begin{itemize}
    \item \textbf{Infraestrutura de Desenvolvimento:}  
    Computador ou notebook com acesso à internet para o aluno e para o professor orientador. O UDF disponibiliza internet para docentes e discentes sem custo. Docente e discente já possuem notebook próprios.
    \item \textbf{Ferramentas de Coleta e Armazenamento:}  
    Acesso à \textit{API do GitHub} e serviços de banco de dados (MongoDB) em servidor local ou em nuvem. API do GitHub tem uma cota pública bastante generosa que atende as necessidades desse projeto. A MongoDB como organização disponibiliza um banco sem custos que é mais do que suficiente para a demanda desse projeto.
    \item \textbf{Ferramentas de Programação e Visualização:}  
    Duas linguagens de programação serão utilizadas, \textit{Python} e \textit{JavaScript} para tratamento e análise dos dados, e bibliotecas ou \textit{frameworks} de \textit{front-end} (\textit{React}, \textit{Vue.js}) para a criação do dashboard. Além disso serão utilizadas duas ferramentas para construção dos softwares (Ambientes de Desenvolvimento Integrados (IDE)), Visual Studio Code e PyCharm, ambas com versões livres para a comunidade.
    \item \textbf{Plataforma de Hospedagem:}  
    O serviço Firebase hosting para publicação do protótipo e posterior disponibilização do dashboard final (frontend), sem custos adicionais. E para o backend será utilizado um servidor Linux de uma empresa parceira do LabTech, sem custos para o projeto.
    \item \textbf{Gestão e Reuniões:}  
    Reuniões quinzenais entre orientador e aluno pesquisador; participação da equipe do segundo semestre de 2025 da Fábrica de Software do UDF como usuários-piloto.
\end{itemize}

\subsection{Cronograma de Execução}
\label{sec:cronograma}

\newcommand{\X}{\textbullet}
\begin{table}[ht]
\centering
\caption{Cronograma de Execução do Projeto (abr/2025 a mar/2026)}
\label{tab:cronograma}
\begin{tabular}{|p{7cm}||c|c|c|c|c|c|c|c|c|c|c|c|}
\hline
\multirow{2}{*}{\textbf{Atividade}} & \multicolumn{12}{c|}{\textbf{Meses (M)}} \\
\cline{2-13}
 & \textbf{1} & \textbf{2} & \textbf{3} & \textbf{4} & \textbf{5} & \textbf{6} & \textbf{7} & \textbf{8} & \textbf{9} & \textbf{10} & \textbf{11} & \textbf{12} \\
\hline
Revisão de Literatura (OE1) & \X & \X &   &   &   &   &   &   &   &   &   &   \\
\hline
Definição de KPIs (OE2)    & \X & \X & \X &   &   &   &   &   &   &   &   &   \\
\hline
Coleta/extração de dados (OE3)  &   &   &   & \X & \X &   &   &   &   &   &   &   \\
\hline
Construção do Dashboard (OE4) &   &   &   &   & \X & \X &   &   &   &   &   &   \\
\hline
Testes-Piloto e Ajustes (OE5) &   &   &   &   &   &   & \X & \X & \X &   &   &   \\
\hline
Análise de Resultados e Relatório &   &   &   &   &   &   &   &   &   & \X & \X &   \\
\hline
Apresentação no Congresso UDF &   &   &   &   &   &   &   &   &   &   &   & \X \\
\hline
\end{tabular}
\end{table}

\begin{itemize}
    \item \textbf{Meses 1--2 (abr/2025 e mai/2025):} Enfoque nos objetivos de revisão de literatura, buscando encontrar paralelos entre o presente trabalho e trabalhos já realizados que tenham uma relação conceitual (OE1); com base na revisão de literatura vamos começar a definição dos KPIs (OE2).
    \item \textbf{Mês 3 (jun/2025):} Concluir a definição das métricas, deixando claro quais serão os dados a serem coletados e os gráficos e tabelas que estarão presentes no dashboard.
    \item \textbf{Meses 4--5 (jul/2025 e ago/2025):} Implementação da camada de extração de dados a partir da API do GitHub (OE3).
    \item \textbf{Meses 5--6 (ago/2025 e set/2025):} Construção do \emph{dashboard} inicial (OE4) com base nas métricas definidas em OE2.
    \item \textbf{Meses 7--9 (out/2025, nov/2025, dez/2025):} Testes-piloto e consolidação, uso real pela Fábrica de Software (OE5). Nessa etapa ira-se-\'{a} entender o uso do dashboard no LabTech, receber feedbacks em relação ao seu uso e fazer algumas adaptações para melhor atender aos usuários.
    \item \textbf{Meses 10--11 (jan/2026 e fev/2026):} Análise de resultados, escrita do relatório final e disponibilização da solução (GPL3).
    \item \textbf{Mês 12 (mar/2026):} Apresentação no Congresso UDF e encerramento do projeto.
\end{itemize}

% ----------------------------------------------------------
% CONSIDERAÇÕES FINAIS
% ----------------------------------------------------------
\section[Considerações Finais]{Considerações Finais}
\addcontentsline{toc}{chapter}{Considerações Finais}

O presente projeto de Iniciação Tecnológica propõe o desenvolvimento de um \textbf{dashboard} para monitoramento de Fábricas de Software Acadêmicas, endereçando a lacuna de padronização na coleta e visualização de métricas de progresso e colaboração. Acredita-se que a solução possa contribuir de forma decisiva para a melhoria da \emph{governança} dos projetos, a detecção precoce de gargalos e o aprimoramento contínuo das equipes, reforçando a formação prática dos estudantes \cite{oliveira2017conduccao, almeida2006liccoes, dos2021experiencia, neto2017relato}.

O cronograma prevê a conclusão do projeto em 12 meses, com atividades que incluem revisão bibliográfica, definição de indicadores, desenvolvimento de um protótipo, construção do dashboard e testes-piloto. Ao final, espera-se disponibilizar os resultados e o código sob \textit{licença GPL3}, bem como apresentar as conclusões no Congresso UDF.

\noindent\textbf{Contribuições esperadas:}
\begin{itemize}
    \item \textbf{Inovação Pedagógica:} Ao consolidar métricas relevantes e fornecer \emph{feedback} mais imediato aos alunos.
    \item \textbf{Pesquisa Aplicada:} Integração de coleta de dados via API do GitHub com visualizações que combinam aspectos técnicos e de colaboração.
    \item \textbf{Expansão Futuras:} Adaptabilidade do \emph{framework} para outras Fábricas de Software em diferentes instituições, favorecendo a pesquisa comparativa.
\end{itemize}

Diante do exposto, identifica-se a oportunidade de conceber e avaliar um \textbf{dashboard} que consolide métricas tanto de ordem técnica (ex.: \emph{commits}, \emph{pull requests}, testes unitários) quanto de ordem colaborativa (ex.: rede de contribuições, frequência e qualidade das interações em \emph{issues}, \emph{code reviews}) \cite{barbosaIndicadoresConjunto2016, hamer2021using}. A adoção de metodologias que sistematizem a coleta e a análise de tais dados configura-se como fator crítico de sucesso para o desenvolvimento sustentável de Fábricas de Software Acadêmicas, oferecendo uma base sólida para tomada de decisões e promoção de melhorias contínuas \cite{fernandes2004fabrica, romanha2016modelo}.

Por meio dessa iniciativa, busca-se não apenas aferir indicadores de produtividade e colaboração, mas também fomentar práticas que incentivem o desenvolvimento de \emph{soft skills} e \emph{hard skills} relevantes ao mercado de TI, aproximando os estudantes de uma realidade profissional cada vez mais exigente e dinâmica \cite{wef:2025, neto2017relato}.

% ----------------------------------------------------------
% Referências Bibliográficas
% ----------------------------------------------------------
\bibliography{abntex2-modelo-references}

% ----------------------------------------------------------
% APÊNDICES (se necessários)
% ----------------------------------------------------------
%\begin{apendicesenv}
%\chapter{Título do Apêndice}
%\lipsum[1]
%\end{apendicesenv}

% ----------------------------------------------------------
% ANEXOS (se necessários)
% ----------------------------------------------------------
%\begin{anexosenv}
%\chapter{Título do Anexo}
%\lipsum[2]
%\end{anexosenv}

\end{document}
